%==============================================================================
%  MATH 231 -- Linear Algebra I (Queens College, Fall 2026)
%  HOMEWORK 3 -- Matrices and Matrix Multiplication
%
%  STATUS: final.  Started 2026-10-12, approved and distributed 2026-10-25.
%  Due Sunday, November 8, 2026, 11:59 PM.
%
%  SCOPE.  This file covers:
%     lectures/matrix-unit/MATH231_matrices_1_linear_systems_and_matrix_notation
%                                        (S2.3, S2.5, S3-S9; lectured Oct 6, 8)
%
%  MATERIAL NOT IN THE LECTURES, defined inside the problem that uses it:
%   - skew-symmetric matrices (B.3).
%
%  WRITTEN ONLY.  There is no programming part.  Any NumPy companion belongs in
%  a separate, optional "Extra Credit 2" assignment, which would refer to this
%  assignment's problems by number.
%
%  NUMBERING.  As in Homeworks 1 and 2: parts are lettered, problems numbered
%  by part.  Bare section numbers refer to the MATRIX unit, whose numbering
%  restarts at S1.  References into the vector unit always say so.
%
%  DELIBERATE OMISSIONS, so that nothing here outruns the lectures:
%   - No Gaussian elimination, no solving of systems beyond reading off
%     coefficients by the column reading (D.3(b)).
%   - No inverses, no determinants.
%
%  CUT AT THE INSTRUCTOR'S REVIEW, 2026-10-25:
%   - the trace problem (was D.6), deleted;
%   - Part E (permutation matrices), deleted;
%   - Parts F (moving shapes with matrices), G (the projection matrix) and
%     H (the Hessian), moved to plan/future_homework/ as catalog Problems
%     4-14.
%
%  NOT INCLUDED: the three problems held back from Homework 2 in
%  plan/future_homework/ (the gradient as an object; the directional
%  derivative; gradient descent and linear approximation).
%
%  NO ANSWER KEY, by design, as in Homeworks 1 and 2.
%
%  The submission path in the "What to submit" callout is homework/hw3/ in
%  the student repo.
%
%  Compile with pdflatex or tectonic (standard TeX Live packages only).
%==============================================================================

\documentclass[11pt]{article}

\usepackage[margin=1in]{geometry}
\usepackage{amsmath,amssymb}
\usepackage[dvipsnames]{xcolor}
\usepackage{enumitem}
\usepackage{booktabs}
\usepackage{array}
\usepackage[hidelinks]{hyperref}
\usepackage{titlesec}
\usepackage{needspace}

%------------------------------------------------------------------ style ----
\titleformat{\section}{\large\bfseries\sffamily\color{RoyalBlue}}
                      {Part \thesection.}{0.6em}{}
\titleformat{\subsection}{\normalsize\bfseries\sffamily}{\thesubsection.}{0.5em}{}
\titlespacing*{\section}{0pt}{16pt}{7pt}

\renewcommand{\thesection}{\Alph{section}}

% Keeps long inline formulas from pushing a line into the margin.
\setlength{\emergencystretch}{3em}

% The problem heading is a line of its own, so that part (a) -- or the opening
% sentence -- always starts below it.  (An amsthm theorem cannot promise this: a
% list that opens the body runs its first item onto the heading line.)  The
% extra air above and below keeps consecutive questions reading as separate
% blocks rather than one continuous column of text.
\makeatletter
\newcounter{prob}[section]
\renewcommand{\theprob}{\thesection.\arabic{prob}}
\newenvironment{prob}[1][]
  {\par\addvspace{20pt plus 6pt minus 4pt}%
   \needspace{4\baselineskip}%
   \refstepcounter{prob}%
   \noindent\textbf{Problem~\theprob}%
   \if\relax\detokenize{#1}\relax\else\ (#1)\fi\textbf{.}%
   \par\nobreak\@afterindentfalse\@afterheading}
  {\par\addvspace{20pt plus 6pt minus 4pt}}
\makeatother

\newcommand{\heads}[1]{\par\smallskip\noindent
  {\small\sffamily\textcolor{BrickRed}{\textbf{Watch out.}}\ #1}\par\smallskip}
\newcommand{\hint}[1]{\par\smallskip\noindent
  {\small\sffamily\textcolor{RoyalBlue}{\textbf{Hint.}}\ #1}\par\smallskip}

% shaded call-out box.  The body is collected into a savebox first, because
% \colorbox needs its whole argument at once.
\definecolor{softbg}{HTML}{F2F6FA}
\newsavebox{\qcbox}
\newenvironment{callout}[1]
  {\par\medskip\needspace{5\baselineskip}%
   \begin{lrbox}{\qcbox}%
   \begin{minipage}{\dimexpr\textwidth-2\fboxsep-4pt}%
   \setlength{\parskip}{0.4em}%
   \textbf{\sffamily\color{RoyalBlue}#1}\par}
  {\end{minipage}\end{lrbox}%
   \begin{center}\colorbox{softbg}{\usebox{\qcbox}}\end{center}\medskip}
%--------------------------------------------------------------- notation ----
\newcommand{\R}{\mathbb{R}}
\newcommand{\vc}[1]{\mathbf{#1}}
\newcommand{\mt}[1]{\mathbf{#1}}          % matrices: bold upper case
\newcommand{\T}{^{\mathsf{T}}}            % transpose
\newcommand{\uv}[1]{\hat{\vc{#1}}}
\newcommand{\one}{\mathbf{1}}
\newcommand{\norm}[1]{\lVert #1\rVert}
\newcommand{\ip}[2]{\langle #1,#2\rangle}
\newcommand{\proj}[2]{\operatorname{proj}_{#1}(#2)}
\DeclareMathOperator{\diagop}{diag}
\DeclareMathOperator{\tr}{tr}
\newcommand{\grad}{\nabla}
\newcommand{\pd}[2]{\frac{\partial #1}{\partial #2}}
\newcommand{\pdd}[3]{\frac{\partial^{2} #1}{\partial #2\,\partial #3}}

%--------------------------------------------------------- fill this in ------
\newcommand{\duedate}{Sunday, November 8, 2026, 11:59 PM}

\title{\vspace{-1.2cm}\textbf{\sffamily MATH 231 -- Linear Algebra I}\\[4pt]
       {\large\sffamily Homework 3 \textbullet\ Matrices and Matrix
        Multiplication}\\[2pt]
       {\normalsize\sffamily Kiely Hall 326 \textbullet\
        Instructor: Kevin Bedoya}}
\author{}
\date{}

\begin{document}
\maketitle
\vspace{-1.6cm}

\noindent\textbf{Due:} \duedate.

\medskip
\begin{center}
\renewcommand{\arraystretch}{1.25}
\begin{tabular}{@{}l l >{\raggedright\arraybackslash}p{7.2cm}@{}}
\toprule
\textbf{Lecture} & \textbf{Sections} & \textbf{Topics} \\
\midrule
Matrix Unit, Part 1 & \S2.3, \S2.5, \S3--\S9 &
  matrices, entries and dimensions; shapes and patterns; adding and scaling;
  the transpose and symmetric matrices; inner and outer products; matrix
  multiplication; matrix--vector products by rows and by columns; the
  trace \\
\bottomrule
\end{tabular}
\end{center}

\noindent\textbf{Grading.} Every problem is worth the same number of points.

\medskip
\noindent\textbf{Section numbers.} A bare section number such as \S8.4 refers
to the matrix unit, \emph{Linear Systems and Matrix Notation}, whose numbering
restarts at $1$. References into the vector unit say so.

\medskip
\noindent\textbf{On the proof problems} (A.4(b), B.1(e), B.2(d), B.3(a), (b),
(d), C.2(a), (c), (d), D.3(c), D.4(d), and D.5(c)). None of them requires a
clever idea. What they require is that you write down what you are assuming, work
entry by entry or from the stated laws, justify each equality by naming the
fact that licenses it, and end with a sentence saying what you have shown. A
chain of equalities with no words is not a proof; neither is a paragraph of
words with no computation. To show two matrices are equal, show that they have
the same shape and that their $(i,j)$ entries agree for every $i$ and $j$.

\medskip
\noindent\textbf{Collaboration} is encouraged --- work at a board together, argue
about the problems. What you submit must be written by you, in your own words.

\begin{callout}{What to submit}
Push your written work for Parts A--D to your own MATH 231 GitHub repository,
as a single PDF named \texttt{hw3.pdf} in a folder named \texttt{homework/hw3/}.
Handwritten and scanned or photographed is completely fine, as long as it is
legible and the pages are in order.

\textbf{PDF is the only accepted format.} A \texttt{.docx}, a \texttt{.jpg}, a
\texttt{.png}, a \texttt{.heic}, or a link to a file hosted somewhere else will
not be graded. Phone scanning apps export PDF directly, and every word processor
can ``Save as PDF''.

\textbf{Type Parts A--D up instead and you earn $+5$ points of extra credit on
this assignment.} Any typesetting system counts --- \LaTeX, Word's equation
editor, or anything else that produces readable mathematics --- and the extra
credit is for the whole of Parts A--D, not for typing up some of it. The file
you submit must still be a PDF. Sketches may be drawn by hand and scanned into
the PDF.

If you have not set up the repository yet, Setup Guide 1 (\S\,``Create your
folders'' and \S\,``Save your work and send it to GitHub'') is the walkthrough.
Only pushed work can be graded.
\end{callout}

\begin{callout}{Label every answer}
Label each answer with the number of the problem it answers --- \texttt{A.3(d)},
\texttt{C.2(b)}, \texttt{D.2(e)} --- and put your answers \textbf{in the same
order as they appear in this assignment}. A grader who cannot tell which
question a piece of work belongs to cannot award it credit, and that is the most
common way a correct answer loses points.

\textbf{You do not have to copy the question out.} This applies to handwritten
and typed submissions alike: the label alone is enough. Write \texttt{A.3(d)},
then your work. Do not reproduce the wording of the problem.
\end{callout}

%==============================================================================
\section{Matrices, entries, and patterns}
%==============================================================================
\emph{Matrix Unit, \S3--\S4.} An $m\times n$ matrix has $m$ rows and $n$
columns, and its entry in row $i$ and column $j$ is $a_{ij}$ --- row first,
column second (\S3.1). $\R^{m\times n}$ is the set of all $m\times n$ matrices
with real entries (\S3.2). A $1\times n$ matrix is a \textbf{row vector} and an
$m\times1$ matrix a \textbf{column vector} (\S3.4). The \textbf{main diagonal}
of a square matrix is $a_{11},a_{22},\dots,a_{nn}$, and the patterns of \S4 ---
square, rectangular, diagonal, identity, upper and lower triangular, banded and
tridiagonal, sparse and dense --- are summarized in the table of \S4.7.

%------------------------------------------------------------------------------
\begin{prob}[Reading a matrix]
Let
\[
  \mt{A}=\begin{bmatrix} 2&-1&0&5\\ 3&4&-2&1\\ 0&7&6&-3 \end{bmatrix},
  \qquad
  \mt{B}=\begin{bmatrix} 1&0\\ -4&9\\ 2&2\\ 8&-5 \end{bmatrix},
  \qquad
  \vc{c}=\begin{bmatrix} 3\\ -1\\ 4 \end{bmatrix},
  \qquad
  \vc{d}=\begin{bmatrix} 6&0&-2&1&5 \end{bmatrix}.
\]
\begin{enumerate}[label=\textbf{(\alph*)},leftmargin=2.6em,itemsep=7pt,topsep=6pt]
  \item Give the dimensions of each, the set $\R^{m\times n}$ it belongs to, and
        its number of entries.
  \item Read off $a_{23}$, $a_{32}$, $a_{14}$, $b_{41}$, $b_{22}$, $c_{31}$ and
        $d_{15}$.
  \item Which of the addresses $a_{41}$, $b_{23}$, $c_{13}$ and $d_{51}$ exist?
        For each one that does not, say why not.
  \item Write row $2$ of $\mt{A}$ as a row vector and column $3$ of $\mt{A}$ as
        a column vector.
  \item Read $\mt{A}$ by rows and by columns, as in \S3.5: how many row vectors
        does it stack, and in which $\R^{k}$ does each one lie? How many column
        vectors, and in which $\R^{k}$?
\end{enumerate}
\end{prob}

%------------------------------------------------------------------------------
\begin{prob}[Building a matrix from a rule]
Write out each matrix entry by entry.
\begin{enumerate}[label=\textbf{(\alph*)},leftmargin=2.6em,itemsep=7pt,topsep=6pt]
  \item $\mt{A}\in\R^{3\times4}$ with $a_{ij}=i-j$.
  \item $\mt{B}\in\R^{4\times4}$ with $b_{ij}=\min(i,j)$.
  \item $\mt{C}\in\R^{4\times4}$ with $c_{ij}=(-1)^{i+j}$.
  \item $\mt{D}\in\R^{4\times4}$ with $d_{ij}=1$ if $\lvert i-j\rvert=1$, and
        $d_{ij}=0$ otherwise.
  \item $\mt{E}\in\R^{4\times4}$ with $e_{ij}=ij$ if $i\le j$, and $e_{ij}=0$
        if $i>j$.
  \item $\mt{F}\in\R^{2\times5}$ with $f_{ij}=10i+j$. What do you notice about
        each entry and its address?
  \item For each of $\mt{B}$, $\mt{C}$, $\mt{D}$ and $\mt{E}$, list every
        pattern from the table of \S4.7 that it has, leaving out sparse and
        dense, and count its nonzero entries.
  \item Give a rule for $a_{ij}$ that produces the identity $\mt{I}_n$, and one
        that produces the $n\times n$ upper triangular matrix whose entries on
        and above the diagonal are all $1$.
\end{enumerate}
\end{prob}

%------------------------------------------------------------------------------
\begin{prob}[Naming the pattern]
For each matrix, list every word that applies to it from: \emph{square,
rectangular, diagonal, identity, upper triangular, lower triangular,
tridiagonal}. Recall that \S4 defines diagonal, identity and triangular
matrices only among square matrices, and that a matrix is tridiagonal when
$a_{ij}=0$ whenever $\lvert i-j\rvert>1$.
\[
  \text{(a) }\begin{bmatrix}0&0&0\\0&0&0\\0&0&0\end{bmatrix}
  \qquad
  \text{(b) }\begin{bmatrix}2&0\\0&2\end{bmatrix}
  \qquad
  \text{(c) }\begin{bmatrix}1&0&0\\0&1&0\end{bmatrix}
  \qquad
  \text{(d) }\begin{bmatrix}4&0&0\\1&0&0\\2&3&0\end{bmatrix}
\]
\[
  \text{(e) }\begin{bmatrix}0&0&1\\0&1&0\\1&0&0\end{bmatrix}
  \qquad
  \text{(f) }\begin{bmatrix}1&2&0&0\\3&4&5&0\\0&6&7&8\\0&0&9&1\end{bmatrix}
  \qquad
  \text{(g) }\begin{bmatrix}5\end{bmatrix}
\]
\begin{enumerate}[label=\textbf{(\alph*)},leftmargin=2.6em,itemsep=7pt,topsep=6pt]
  \setcounter{enumi}{7}
  \item True or false? Justify each answer from the definitions in \S4, and give
        a counterexample for each false statement.
        \begin{enumerate}[label=(\roman*),leftmargin=2.2em,itemsep=3pt,topsep=3pt]
          \item Every diagonal matrix is upper triangular.
          \item Every identity matrix is tridiagonal.
          \item An upper triangular matrix whose diagonal entries are all $0$ is
                the zero matrix.
          \item If $\mt{A}$ is lower triangular, then so is $c\mt{A}$ for every
                scalar $c$.
        \end{enumerate}
\end{enumerate}
\end{prob}

%------------------------------------------------------------------------------
\begin{prob}[The diagonal of any matrix, and counting entries]
The main diagonal makes sense for a rectangular matrix too: for
$\mt{A}\in\R^{m\times n}$ it consists of the entries $a_{11},a_{22},a_{33},\dots$,
continuing for as long as both row $i$ and column $i$ exist.
\begin{enumerate}[label=\textbf{(\alph*)},leftmargin=2.6em,itemsep=7pt,topsep=6pt]
  \item List the main diagonal of each matrix:
        \[
          \begin{bmatrix}1&2&3&4&5\\6&7&8&9&10\end{bmatrix},
          \qquad
          \begin{bmatrix}1&2&3\\4&5&6\\7&8&9\\10&11&12\end{bmatrix},
          \qquad
          \begin{bmatrix}3&1&4&1&5&9\end{bmatrix},
          \qquad
          \begin{bmatrix}2\\7\\1\\8\end{bmatrix}.
        \]
  \item Show that an $m\times n$ matrix has exactly $\min(m,n)$ entries on its
        main diagonal.
  \item In an $n\times n$ matrix, count the entries that lie on the main
        diagonal, strictly above it, strictly below it, and on or above it.
        Give each count as a formula in $n$, and check your formulas against a
        drawing of a $4\times4$ grid.
        \hint{The entries strictly above the diagonal in row $i$ are those in
        columns $i+1,\dots,n$. Add up over the rows, and use the sum formula for
        $1+2+\cdots+k$.}
  \item The \emph{largest} number of entries that can be nonzero in an
        $n\times n$ matrix that is diagonal is $n$. Find the corresponding
        largest number for an upper triangular matrix and for a tridiagonal
        matrix.
  \item For a tridiagonal $n\times n$ matrix, compute the fraction of its $n^{2}$
        entries that can be nonzero, for $n=4$, $n=10$ and $n=1000$. What
        happens to this fraction as $n\to\infty$?
  \item A $1000\times1000$ matrix has $4000$ nonzero entries. What fraction of
        its entries are zero? Storing the full grid takes one number per entry.
        Storing only the nonzero entries takes three numbers per nonzero entry:
        its row, its column, and its value. How many numbers does each way take?
        By the working test of \S4.6, is this matrix sparse?
\end{enumerate}
\end{prob}

%==============================================================================
\section{Adding, scaling, and the transpose}
%==============================================================================
\emph{Matrix Unit, \S5--\S6.} Matrices of the same shape are added and
subtracted entry by entry, $(\mt{A}\pm\mt{B})_{ij}=a_{ij}\pm b_{ij}$, and
scaled entry by entry, $(c\mt{A})_{ij}=ca_{ij}$. The laws of addition and
scaling are \textbf{A1}--\textbf{A4} and \textbf{S1}--\textbf{S4}, introduced in
\S5.4 (``The algebraic laws''). The transpose of $\mt{A}\in\R^{m\times n}$ is
$\mt{A}\T\in\R^{n\times m}$ with $(\mt{A}\T)_{ij}=a_{ji}$ (\S6.1), and its laws
are
\[
  \textbf{T1}\ \ (\mt{A}\T)\T=\mt{A},
  \qquad
  \textbf{T2}\ \ (\mt{A}+\mt{B})\T=\mt{A}\T+\mt{B}\T,
  \qquad
  \textbf{T3}\ \ (c\mt{A})\T=c\,\mt{A}\T
\]
(\S6.4). A matrix is \textbf{symmetric} if $\mt{A}\T=\mt{A}$ (\S6.5).

%------------------------------------------------------------------------------
\begin{prob}[Matrix arithmetic]
Let
\[
  \mt{A}=\begin{bmatrix}1&-2&0\\3&4&-1\end{bmatrix},
  \qquad
  \mt{B}=\begin{bmatrix}2&0&5\\-1&3&2\end{bmatrix},
  \qquad
  \mt{C}=\begin{bmatrix}1&4\\-2&0\\3&1\end{bmatrix}.
\]
\begin{enumerate}[label=\textbf{(\alph*)},leftmargin=2.6em,itemsep=7pt,topsep=6pt]
  \item Compute $\mt{A}+\mt{B}$, $\mt{A}-\mt{B}$, $-\mt{B}$ and
        $3\mt{A}-2\mt{B}$.
  \item Which of $\mt{A}+\mt{C}$, $\mt{A}+\mt{C}\T$, $\mt{A}\T+\mt{C}$ and
        $\mt{B}-\mt{C}\T$ are defined? Compute each one that is, and for each one
        that is not, say why.
  \item Compute $\mt{A}\T$, $\mt{B}\T$, $(\mt{A}+\mt{B})\T$ and
        $\mt{A}\T+\mt{B}\T$, and confirm \textbf{T2} for this pair.
  \item Find the matrix $\mt{X}$ with $3\mt{X}-\mt{B}=\mt{X}+\mt{A}$. Solve for
        $\mt{X}$ first, naming the law from \S5.4 used at each step, and only
        then substitute.
  \item Find scalars $\alpha$ and $\beta$ with
        \[
          \alpha\begin{bmatrix}1&0\\0&1\end{bmatrix}
          +\beta\begin{bmatrix}0&1\\1&0\end{bmatrix}
          =\begin{bmatrix}3&-2\\-2&3\end{bmatrix}.
        \]
        Then show that $\begin{bmatrix}1&2\\3&4\end{bmatrix}$ cannot be written
        in this form for any $\alpha$ and $\beta$.
\end{enumerate}
\end{prob}

%------------------------------------------------------------------------------
\begin{prob}[Transposes and symmetric matrices]
\begin{enumerate}[label=\textbf{(\alph*)},leftmargin=2.6em,itemsep=7pt,topsep=6pt]
  \item Write down the transpose of each matrix, with its shape:
        \[
          \begin{bmatrix}1&2&3&4\end{bmatrix},
          \qquad
          \begin{bmatrix}5\\-1\\2\end{bmatrix},
          \qquad
          \begin{bmatrix}0&1\\2&3\\4&5\end{bmatrix},
          \qquad
          \begin{bmatrix}1&2&3\\2&5&4\\3&4&6\end{bmatrix}.
        \]
        Which of the four is symmetric?
  \item Find every $x$, $y$ and $z$ for which
        $\begin{bmatrix}1&x&3\\4&2&y\\z&-5&0\end{bmatrix}$ is symmetric.
  \item Fill in the entries marked $\ast$ so that the matrix is symmetric:
        $\begin{bmatrix}2&\ast&\ast\\7&-1&\ast\\0&3&4\end{bmatrix}$.
  \item Prove that if $\mt{A}$ and $\mt{B}$ are symmetric $n\times n$ matrices
        and $c$ is a scalar, then $\mt{A}+\mt{B}$ and $c\mt{A}$ are symmetric.
        Name the law from \S6.4 used at each step.
\end{enumerate}
\end{prob}

%------------------------------------------------------------------------------
\begin{prob}[Splitting a square matrix into symmetric and skew-symmetric parts]
A square matrix $\mt{K}$ is \textbf{skew-symmetric} if $\mt{K}\T=-\mt{K}$, that
is, if $k_{ji}=-k_{ij}$ for every $i$ and $j$. In (a), (b) and (d), work from the
definitions and the laws \textbf{T1}--\textbf{T3}, \textbf{A1}--\textbf{A4} and
\textbf{S1}--\textbf{S4}, naming each one you use.
\begin{enumerate}[label=\textbf{(\alph*)},leftmargin=2.6em,itemsep=7pt,topsep=6pt]
  \item Prove that every diagonal entry of a skew-symmetric matrix is $0$.
  \item Prove that for every square matrix $\mt{A}$, the matrix
        $\mt{A}+\mt{A}\T$ is symmetric and the matrix $\mt{A}-\mt{A}\T$ is
        skew-symmetric.
  \item It follows that every square matrix is the sum of a symmetric matrix and
        a skew-symmetric one:
        \[
          \mt{A}
          =\underbrace{\tfrac12\bigl(\mt{A}+\mt{A}\T\bigr)}_{\text{symmetric}}
          +\underbrace{\tfrac12\bigl(\mt{A}-\mt{A}\T\bigr)}_{\text{skew-symmetric}} .
        \]
        Compute both parts for
        $\mt{A}=\begin{bmatrix}1&4&-2\\0&3&5\\6&-1&2\end{bmatrix}$, and check
        that they add up to $\mt{A}$.
  \item Prove that the only matrix that is both symmetric and skew-symmetric is
        the zero matrix.
\end{enumerate}
\end{prob}

%==============================================================================
\section{Inner and outer products}
%==============================================================================
\emph{Matrix Unit, \S7.} From \S6.3 on, a vector $\vc{v}\in\R^{n}$ is a column,
an $n\times1$ matrix, and $\vc{v}\T$ is the same vector written as a row. The
bracket notation $[\,v_1,\dots,v_n\,]$ is still used inline, and means the
column. For $\vc{u},\vc{v}\in\R^{n}$ the \textbf{inner product}
$\vc{u}\T\vc{v}=\sum_{k=1}^{n}u_kv_k$ is the dot product, a $1\times1$ matrix
read as a scalar (\S7.1). For $\vc{u}\in\R^{m}$ and $\vc{v}\in\R^{n}$ the
\textbf{outer product} $\vc{u}\vc{v}\T$ is the $m\times n$ matrix whose $(i,j)$
entry is $u_iv_j$ (\S7.2). And $\vc{v}\T\vc{v}=\norm{\vc{v}}^{2}$ (\S7.3).

%------------------------------------------------------------------------------
\begin{prob}[Computing inner and outer products]
Let $\vc{u}=[\,1,\,-2,\,3\,]$, $\vc{v}=[\,4,\,0,\,-1\,]$ and
$\vc{w}=[\,2,\,5\,]$. Before computing each product, state its shape.
\begin{enumerate}[label=\textbf{(\alph*)},leftmargin=2.6em,itemsep=7pt,topsep=6pt]
  \item $\vc{u}\T\vc{v}$, $\vc{v}\T\vc{u}$, $\vc{u}\T\vc{u}$ and
        $\vc{v}\T\vc{v}$. Use the last two to find $\norm{\vc{u}}$ and
        $\norm{\vc{v}}$.
  \item $\vc{u}\vc{v}\T$ and $\vc{v}\vc{u}\T$. Check that
        $(\vc{u}\vc{v}\T)\T=\vc{v}\vc{u}\T$.
  \item $\vc{u}\vc{w}\T$ and $\vc{w}\vc{u}\T$. Is $\vc{u}\T\vc{w}$ defined?
  \item Use $\norm{\vc{x}}=\sqrt{\vc{x}\T\vc{x}}$ to compute the norms of
        $[\,2,\,-1,\,2,\,4\,]\in\R^{4}$ and of the all-ones vector
        $\one\in\R^{9}$.
\end{enumerate}
\end{prob}

%------------------------------------------------------------------------------
\begin{prob}[What an outer product looks like]
\begin{enumerate}[label=\textbf{(\alph*)},leftmargin=2.6em,itemsep=7pt,topsep=6pt]
  \item Let $\vc{u}\in\R^{m}$ and $\vc{v}\in\R^{n}$. Prove that row $i$ of
        $\vc{u}\vc{v}\T$ is $u_i\vc{v}\T$, and that column $j$ of
        $\vc{u}\vc{v}\T$ is $v_j\vc{u}$. So every row is a multiple of the same
        row vector, and every column is a multiple of the same column vector.
  \item Each matrix below is an outer product $\vc{u}\vc{v}\T$. Find a pair
        $\vc{u}$, $\vc{v}$ for each, and check it by multiplying out.
        \[
          \begin{bmatrix}2&4&6\\3&6&9\end{bmatrix},
          \qquad
          \begin{bmatrix}1&-1\\-2&2\\3&-3\end{bmatrix},
          \qquad
          \begin{bmatrix}0&0&0\\5&10&-5\\-1&-2&1\end{bmatrix},
          \qquad
          \begin{bmatrix}6&-4&2&8\end{bmatrix}.
        \]
        Is the pair unique? If not, describe every other pair that works for the
        first matrix.
  \item Prove that $\begin{bmatrix}1&2\\3&4\end{bmatrix}$ is \emph{not} an outer
        product $\vc{u}\vc{v}\T$ of any two vectors $\vc{u},\vc{v}\in\R^{2}$.
  \item Prove that $\tr(\vc{u}\vc{v}\T)=\vc{v}\T\vc{u}$ for all
        $\vc{u},\vc{v}\in\R^{n}$, where the \textbf{trace} of a square matrix is
        the sum of its diagonal entries (\S9.1). Check it on $\vc{u}$ and
        $\vc{v}$ from Problem C.1.
  \item Write $\begin{bmatrix}1&2\\3&4\end{bmatrix}$ as the sum of two outer
        products.
        \hint{Try $\vc{e}_1$ and $\vc{e}_2$ as the column vectors.}
\end{enumerate}
\end{prob}

%==============================================================================
\section{Matrix multiplication}
%==============================================================================
\emph{Matrix Unit, \S2.3, \S2.5 and \S8.} For $\mt{A}\in\R^{m\times p}$ and
$\mt{B}\in\R^{p\times n}$, the product $\mt{A}\mt{B}\in\R^{m\times n}$ has
$(i,j)$ entry
\[
  (\mt{A}\mt{B})_{ij}=\sum_{k=1}^{p}a_{ik}b_{kj},
\]
the dot product of row $i$ of $\mt{A}$ with column $j$ of $\mt{B}$ (\S8.1). It
is defined exactly when the number of columns of $\mt{A}$ equals the number of
rows of $\mt{B}$ (\S8.2). The laws are \textbf{M1}--\textbf{M6} of \S8.4
(``The algebraic laws''): associativity, the two distributive laws, scalars
moving freely, $\mt{I}\mt{A}=\mt{A}\mt{I}=\mt{A}$, and
$(\mt{A}\mt{B})\T=\mt{B}\T\mt{A}\T$. There is no commutative law. A
matrix--vector product $\mt{A}\vc{x}$ is the case $n=1$ (\S8.5), and it can be
read two ways (\S2.5): by rows, its $i$-th entry is the dot product of row $i$
of $\mt{A}$ with $\vc{x}$; by columns, it is the linear combination
$x_1\vc{a}_1+\cdots+x_p\vc{a}_p$ of the columns $\vc{a}_1,\dots,\vc{a}_p$ of
$\mt{A}$. (\S2.5 states this for a square matrix; the proof given there works
unchanged for any $\mt{A}\in\R^{m\times p}$.) Write $\vc{e}_i$ for
the standard basis vector with a $1$ in slot $i$ and $0$ elsewhere.

%------------------------------------------------------------------------------
\begin{prob}[Shapes first]
Let $\mt{A}\in\R^{2\times3}$, $\mt{B}\in\R^{3\times3}$, $\mt{C}\in\R^{3\times1}$,
$\mt{D}\in\R^{1\times2}$ and $\mt{E}\in\R^{2\times2}$. For each product, say
whether it is defined, and if it is, give its shape. Set your answers out as a
table.
\[
  \mt{A}\mt{B},\ \ \mt{B}\mt{A},\ \ \mt{A}\mt{C},\ \ \mt{C}\mt{A},\ \
  \mt{B}\mt{C},\ \ \mt{C}\mt{B},\ \ \mt{D}\mt{A},\ \ \mt{A}\mt{D},\ \
  \mt{D}\mt{E},\ \ \mt{E}\mt{D},\ \ \mt{E}\mt{A},\ \ \mt{A}\mt{E},
\]
\[
  \mt{C}\T\mt{B},\ \ \mt{C}\mt{C}\T,\ \ \mt{C}\T\mt{C},\ \ \mt{A}\T\mt{E},\ \
  \mt{A}\mt{A}\T,\ \ \mt{A}\T\mt{A},\ \ \mt{D}\mt{A}\mt{C},\ \
  \mt{E}\mt{A}\mt{B}\mt{C}.
\]
\end{prob}

%------------------------------------------------------------------------------
\begin{prob}[Multiplying matrices]
\begin{enumerate}[label=\textbf{(\alph*)},leftmargin=2.6em,itemsep=7pt,topsep=6pt]
  \item Let $\mt{E}=\begin{bmatrix}2&-1\\3&4\end{bmatrix}$ and
        $\mt{F}=\begin{bmatrix}1&5\\-2&0\end{bmatrix}$. Compute $\mt{E}\mt{F}$
        and $\mt{F}\mt{E}$. Are they equal?
  \item Let $\mt{A}=\begin{bmatrix}1&0&2\\-1&3&1\end{bmatrix}$ and
        $\mt{B}=\begin{bmatrix}3&1\\2&1\\1&0\end{bmatrix}$. Compute
        $\mt{A}\mt{B}$ and $\mt{B}\mt{A}$.
  \item Let $\mt{G}=\begin{bmatrix}1&2&0\\0&1&-1\\3&0&2\end{bmatrix}$ and
        $\mt{K}=\begin{bmatrix}2&0&1\\1&-1&0\\0&3&1\end{bmatrix}$. Compute
        $\mt{G}\mt{K}$ and $\mt{K}\mt{G}$.
  \item With $\mt{A}$ and $\mt{B}$ from (b), compute $\mt{B}\T\mt{A}\T$ and
        confirm \textbf{M6}. Then compute $\mt{A}\T\mt{B}\T$, and say which
        matrix from (b) it is the transpose of.
  \item Compute
        $\begin{bmatrix}1&2&3\end{bmatrix}\begin{bmatrix}4\\5\\6\end{bmatrix}$
        and
        $\begin{bmatrix}4\\5\\6\end{bmatrix}\begin{bmatrix}1&2&3\end{bmatrix}$,
        and name each as an inner or an outer product.
\end{enumerate}
\end{prob}

%------------------------------------------------------------------------------
\begin{prob}[Matrix--vector products, two ways]
\begin{enumerate}[label=\textbf{(\alph*)},leftmargin=2.6em,itemsep=7pt,topsep=6pt]
  \item Compute each product twice, by rows and by columns, and check that the
        two answers agree:
        \[
          \begin{bmatrix}2&-1&0\\1&3&2\end{bmatrix}
          \begin{bmatrix}3\\1\\-2\end{bmatrix},
          \qquad
          \begin{bmatrix}1&4\\-2&0\\3&1\end{bmatrix}
          \begin{bmatrix}2\\-1\end{bmatrix},
          \qquad
          \begin{bmatrix}1&2&3\\4&5&6\\7&8&9\end{bmatrix}
          \begin{bmatrix}1\\-2\\1\end{bmatrix}.
        \]
        The third product is the zero vector. Write out what the column reading
        then says about the three columns of that matrix. Are the columns
        linearly independent? (Use the definition in the table of the vector
        unit's \S24.2.)
  \item Let $\mt{A}=\begin{bmatrix}1&0&2\\3&-1&1\\0&4&5\end{bmatrix}$, with
        columns $\vc{a}_1,\vc{a}_2,\vc{a}_3$. Use the column reading, and no
        equation solving, to find a vector $\vc{x}$ with each property:
        $\mt{A}\vc{x}=\vc{a}_2$; \ $\mt{A}\vc{x}=3\vc{a}_1-\vc{a}_3$; \
        $\mt{A}\vc{x}=\vc{a}_1+\vc{a}_2+\vc{a}_3$. Check the second one by
        multiplying by rows.
  \item Let $\mt{A}\in\R^{m\times n}$. Prove that $\mt{A}\vc{e}_j$ is column $j$
        of $\mt{A}$, and that $\vc{e}_i\T\mt{A}\vc{e}_j=a_{ij}$. (Here
        $\vc{e}_j\in\R^{n}$ and $\vc{e}_i\in\R^{m}$.) Use the second fact to read
        off $\vc{e}_2\T\mt{A}\vc{e}_3$ for the $\mt{A}$ of (b).
  \item Find the $2\times2$ matrix $\mt{M}$ with
        $\mt{M}\begin{bmatrix}1\\0\end{bmatrix}=\begin{bmatrix}3\\-1\end{bmatrix}$
        and
        $\mt{M}\begin{bmatrix}0\\1\end{bmatrix}=\begin{bmatrix}2\\5\end{bmatrix}$.
        Then compute $\mt{M}\begin{bmatrix}4\\-2\end{bmatrix}$ by the column
        reading.
\end{enumerate}
\end{prob}

%------------------------------------------------------------------------------
\begin{prob}[Matrix arithmetic is not number arithmetic]
\begin{enumerate}[label=\textbf{(\alph*)},leftmargin=2.6em,itemsep=7pt,topsep=6pt]
  \item Let $\mt{A}=\begin{bmatrix}1&2\\2&4\end{bmatrix}$ and
        $\mt{B}=\begin{bmatrix}2&-4\\-1&2\end{bmatrix}$. Compute $\mt{A}\mt{B}$
        and $\mt{B}\mt{A}$. For real numbers, $ab=0$ forces $a=0$ or $b=0$. What
        does (a) show about matrices?
  \item Let $\mt{N}=\begin{bmatrix}0&1\\0&0\end{bmatrix}$ and
        $\mt{N}_3=\begin{bmatrix}0&1&0\\0&0&1\\0&0&0\end{bmatrix}$. Compute
        $\mt{N}^{2}$, then $\mt{N}_3^{2}$ and $\mt{N}_3^{3}$. (Here
        $\mt{A}^{2}=\mt{A}\mt{A}$ and $\mt{A}^{3}=\mt{A}\mt{A}\mt{A}$.)
  \item Use (a) to give matrices with $\mt{A}\mt{B}=\mt{A}\mt{C}$, $\mt{A}\neq\mt{0}$
        and $\mt{B}\neq\mt{C}$. So a nonzero matrix cannot in general be
        cancelled.
  \item Let $\mt{A}$ and $\mt{B}$ be $n\times n$. Prove that
        \[
          (\mt{A}+\mt{B})^{2}=\mt{A}^{2}+\mt{A}\mt{B}+\mt{B}\mt{A}+\mt{B}^{2},
        \]
        naming the law of \S8.4 used at each step, and deduce that
        $(\mt{A}+\mt{B})^{2}=\mt{A}^{2}+2\mt{A}\mt{B}+\mt{B}^{2}$ holds if and
        only if $\mt{A}\mt{B}=\mt{B}\mt{A}$. Then compute both sides of the
        second formula for $\mt{E}$ and $\mt{F}$ from Problem D.2(a).
\end{enumerate}
\end{prob}

%------------------------------------------------------------------------------
\begin{prob}[Products of patterned matrices]
For numbers $d_1,d_2,\dots,d_n$, the notation $\diagop(d_1,d_2,\dots,d_n)$
means the $n\times n$ diagonal matrix (\S4.2) whose diagonal entries are
$d_1,d_2,\dots,d_n$, in that order, and whose other entries are all $0$:
\[
  \diagop(d_1,d_2,\dots,d_n)=
  \begin{bmatrix}
    d_1    & 0      & \cdots & 0\\
    0      & d_2    & \cdots & 0\\
    \vdots & \vdots & \ddots & \vdots\\
    0      & 0      & \cdots & d_n
  \end{bmatrix}.
\]
\begin{enumerate}[label=\textbf{(\alph*)},leftmargin=2.6em,itemsep=7pt,topsep=6pt]
  \item Compute $\diagop(2,-1,3)\,\diagop(4,5,-2)$ and
        $\diagop(4,5,-2)\,\diagop(2,-1,3)$.
  \item Let $\mt{D}=\diagop(2,-1,3)$ and
        $\mt{A}=\begin{bmatrix}1&2&3\\4&5&6\\7&8&9\end{bmatrix}$. Compute
        $\mt{D}\mt{A}$ and $\mt{A}\mt{D}$. In each product, what happened to the
        rows or columns of $\mt{A}$?
  \item Let $\mt{U}=\begin{bmatrix}1&2&3\\0&4&5\\0&0&6\end{bmatrix}$ and
        $\mt{V}=\begin{bmatrix}2&0&1\\0&1&-1\\0&0&3\end{bmatrix}$. Compute
        $\mt{U}\mt{V}$, and observe that it is upper triangular and that its
        diagonal entries are the products $u_{ii}v_{ii}$. Then prove both facts
        for any two upper triangular $n\times n$ matrices.
        \hint{For $i>j$, look at each term $u_{ik}v_{kj}$ of
        $(\mt{U}\mt{V})_{ij}=\sum_{k}u_{ik}v_{kj}$. If $k<i$, which factor is
        zero? If $k\ge i$, which factor is zero?}
  \item Let $\mt{S}=\begin{bmatrix}1&1\\0&1\end{bmatrix}$. Compute $\mt{S}^{2}$,
        $\mt{S}^{3}$ and $\mt{S}^{4}$, guess a formula for $\mt{S}^{n}$, and
        check that your formula satisfies $\mt{S}\,\mt{S}^{n}=\mt{S}^{n+1}$.
  \item Give a formula for $\diagop(2,3)^{n}$.
\end{enumerate}
\end{prob}

\bigskip
\noindent\rule{\textwidth}{0.4pt}

\smallskip
\noindent{\small Questions about any problem here belong in the class Discord,
where everyone can see the answer. Questions that are specific to your own work
--- ``is my proof of D.5(c) valid?'' --- are better brought to office hours,
Thursdays 3:00--4:00 PM in Kiely Hall 508. Please email ahead.}

\end{document}
